%=======================================================================
\chapter{État de l'art}

\section{Introduction}

L'analyse automatique des images médicales constitue aujourd'hui l'un des domaines les plus
actifs de la recherche en intelligence artificielle appliquée à la santé. Grâce aux progrès de
l'imagerie médicale et à l'augmentation de la quantité de données disponibles, les systèmes
d'aide au diagnostic assistés par ordinateur (\emph{Computer-Aided Diagnosis} --- CAD) sont
devenus des outils prometteurs pour améliorer la précision et la rapidité des décisions
cliniques.

Parmi les différentes tâches d'analyse d'images médicales, la classification occupe une place
centrale. Elle consiste à attribuer automatiquement une image à une catégorie prédéfinie en
fonction des caractéristiques qu'elle contient. Dans le domaine de la neuro-oncologie, cette
tâche vise notamment à distinguer différents types de tumeurs cérébrales (gliome, méningiome,
tumeur pituitaire, ou absence de tumeur) à partir d'images d'Imagerie par Résonance
Magnétique (IRM).

L'IRM est considérée comme la modalité d'imagerie de référence pour l'étude du cerveau grâce à
sa capacité à fournir des informations anatomiques détaillées sans exposition aux rayonnements
ionisants. Elle permet de visualiser la localisation, la taille et certaines propriétés des
tumeurs cérébrales, facilitant ainsi leur diagnostic et leur suivi clinique. Cependant,
l'interprétation manuelle des examens IRM demeure une tâche complexe, chronophage et sujette à
une variabilité inter- et intra-observateur susceptible d'affecter la reproductibilité du
diagnostic.

Afin d'améliorer la fiabilité des diagnostics, de nombreuses méthodes d'intelligence
artificielle ont été proposées pour automatiser l'analyse des images médicales. Ces approches
reposent principalement sur deux paradigmes~: le \emph{Machine Learning} traditionnel, fondé
sur l'extraction de caractéristiques explicites suivie d'un classifieur, et le \emph{Deep
Learning}, capable d'apprendre automatiquement des représentations à partir des données brutes.
Ce chapitre présente les fondements de ces deux familles d'approches, discute l'opposition
entre caractéristiques manuelles et caractéristiques profondes, puis dresse une synthèse
critique des travaux récents consacrés à la classification des tumeurs cérébrales.

% -----------------------------------------------------------------------
\section{Machine Learning traditionnel}

Les méthodes de Machine Learning traditionnel reposent généralement sur deux étapes
principales~: l'extraction de caractéristiques descriptives à partir des images, puis
l'apprentissage d'un modèle de classification à partir de ces caractéristiques. Cette
approche présente l'avantage de produire des représentations interprétables tout en
nécessitant moins de données et de ressources de calcul que les méthodes profondes.

\subsection{Support Vector Machine (SVM)}

Le Support Vector Machine (SVM) est l'un des algorithmes les plus utilisés en classification
d'images médicales. Son principe consiste à rechercher l'hyperplan optimal séparant les
différentes classes tout en maximisant la marge entre celles-ci.

Grâce à l'utilisation de fonctions noyaux (\emph{kernels}), le SVM est capable de modéliser
des frontières de décision complexes dans des espaces de grande dimension. Il est
particulièrement adapté aux jeux de données comportant un nombre limité d'exemples mais un
grand nombre de caractéristiques, situation fréquemment rencontrée dans les applications
médicales, ce qui explique sa présence récurrente comme classifieur final dans les approches
hybrides récentes.

\subsection{K-Nearest Neighbors (KNN)}

L'algorithme K-Nearest Neighbors (KNN) repose sur un principe simple~: un nouvel échantillon
est affecté à la classe majoritaire parmi ses $k$ voisins les plus proches dans l'espace des
caractéristiques.

Cette méthode ne nécessite pas de phase d'apprentissage complexe et constitue souvent une
référence efficace pour les problèmes de classification. Toutefois, ses performances peuvent
être dégradées par la présence de caractéristiques redondantes ou par la malédiction de la
dimensionnalité, ce qui rend souvent nécessaire une étape de réduction de dimension (par
exemple par Analyse en Composantes Principales).

\subsection{Random Forest}

Random Forest est une méthode d'apprentissage par ensemble basée sur la combinaison de
plusieurs arbres de décision construits sur des sous-échantillons aléatoires des données.

La prédiction finale est obtenue par vote majoritaire des arbres individuels. Cette approche
permet généralement d'améliorer la robustesse du modèle et de réduire les risques de
surapprentissage tout en offrant des mécanismes d'évaluation de l'importance des
caractéristiques.

\subsection{Extra Trees}

L'algorithme Extra Trees (\emph{Extremely Randomized Trees}) constitue une variante de
Random Forest introduisant davantage d'aléatoire lors de la construction des arbres.

Contrairement à Random Forest, les seuils de séparation sont sélectionnés de manière
aléatoire, ce qui contribue à réduire davantage la variance du modèle et à améliorer sa
capacité de généralisation dans certains contextes.

\subsection{Régression Logistique}

La Régression Logistique (\emph{Logistic Regression}) est une méthode statistique supervisée
largement utilisée pour les tâches de classification.

Bien que relativement simple, elle offre une bonne interprétabilité des résultats et
constitue souvent une référence de base dans les études comparatives. Son efficacité dépend
toutefois de la capacité des caractéristiques à séparer linéairement les différentes classes.

% -----------------------------------------------------------------------
\section{Deep Learning pour l'imagerie médicale}

L'émergence du Deep Learning a profondément transformé le domaine de l'analyse d'images
médicales. Contrairement aux approches traditionnelles, les modèles profonds apprennent
automatiquement les représentations pertinentes directement à partir des pixels de l'image,
supprimant ainsi l'étape d'ingénierie manuelle des caractéristiques.

Cette capacité à extraire des caractéristiques hiérarchiques a permis d'obtenir des
performances remarquables dans de nombreuses applications de détection, segmentation et
classification médicale.

\subsection{Convolutional Neural Networks (CNN)}

Les réseaux de neurones convolutifs (CNN) constituent l'architecture de référence pour le
traitement des images.

Les couches convolutives permettent d'extraire automatiquement des motifs visuels tels que les
contours, les textures ou les structures complexes présentes dans les images médicales. Les
couches profondes apprennent progressivement des représentations de plus en plus abstraites
favorisant la discrimination entre les différentes classes.

\subsection{VGG}

VGG est une architecture profonde caractérisée par l'utilisation répétée de petites
convolutions de taille $3\times3$.

Sa simplicité architecturale et ses bonnes performances ont contribué à sa popularité dans
les applications médicales. Les versions pré-entraînées sur de grands jeux de données
peuvent être adaptées aux tâches médicales grâce à des techniques de transfert
d'apprentissage. Plusieurs travaux rapportent des précisions supérieures à 98\,\% avec des
variantes VGG-16 et VGG-19 en classification binaire de tumeurs cérébrales.

\subsection{ResNet}

ResNet (\emph{Residual Network}) introduit le concept de connexions résiduelles permettant de
faciliter l'apprentissage de réseaux très profonds.

Cette innovation réduit les problèmes de disparition du gradient et améliore
significativement les performances sur les tâches complexes de classification. Les
architectures ResNet figurent aujourd'hui parmi les modèles les plus utilisés dans les
travaux récents portant sur la classification des tumeurs cérébrales, souvent comme extracteur
de caractéristiques au sein d'approches hybrides ou de transfert d'apprentissage.

\subsection{Architectures récentes~: EfficientNet et Vision Transformers}

Au-delà des architectures classiques, les travaux les plus récents s'appuient sur des modèles
plus efficaces ou plus expressifs. EfficientNet propose un équilibrage systématique de la
profondeur, de la largeur et de la résolution du réseau, permettant d'atteindre de hautes
performances avec un nombre réduit de paramètres. Parallèlement, les Vision Transformers (ViT)
et leurs variantes hiérarchiques telles que le Swin Transformer exploitent des mécanismes
d'attention pour capturer des dépendances à longue portée dans l'image, et atteignent
actuellement les meilleures performances rapportées sur les jeux de données de référence.

% -----------------------------------------------------------------------
\section{Handcrafted Features vs Deep Features}

Les approches traditionnelles et profondes présentent chacune des avantages et des
limitations. Le choix entre ces deux paradigmes dépend généralement du volume de données
disponible, des ressources de calcul et des objectifs d'interprétabilité.

\subsection{Familles de caractéristiques manuelles}

Les caractéristiques conçues manuellement (\emph{Handcrafted Features}) sont traditionnellement
regroupées en trois catégories complémentaires~:

\begin{itemize}[leftmargin=1.4em,itemsep=2pt]
  \item les caractéristiques \textbf{d'intensité}, qui décrivent la distribution locale des
        niveaux de gris (moyenne, écart-type, asymétrie, aplatissement, mesures
        histogrammiques)~;
  \item les caractéristiques de \textbf{texture}, qui décrivent l'arrangement spatial des
        intensités~: matrices de cooccurrence (GLCM), motifs binaires locaux (LBP), filtres
        de Gabor et transformées en ondelettes (DWT)~;
  \item les caractéristiques de \textbf{forme}, qui capturent les propriétés géométriques de
        la région tumorale (aire, périmètre, compacité, moments), auxquelles s'ajoutent les
        histogrammes de gradients orientés (HOG) décrivant les contours.
\end{itemize}

Dans ce projet, cinq familles de descripteurs sont étudiées~: Statistiques, GLCM, LBP, DWT et
HOG.

\subsection{Avantages des caractéristiques manuelles}

Les caractéristiques manuelles offrent plusieurs avantages~:

\begin{itemize}[leftmargin=1.4em,itemsep=2pt]
  \item forte interprétabilité des informations extraites~;
  \item besoins réduits en données d'entraînement~;
  \item coût computationnel relativement faible~;
  \item facilité d'analyse et de validation scientifique~;
  \item compatibilité avec de nombreux algorithmes de Machine Learning.
\end{itemize}

\subsection{Limites des caractéristiques manuelles}

Malgré leurs avantages, les caractéristiques manuelles présentent certaines limites. Leur
qualité dépend fortement de l'expertise du concepteur et des choix effectués lors de
l'extraction. Elles peuvent également être moins capables de capturer des représentations
complexes ou de haut niveau présentes dans les données médicales.

Par ailleurs, aucun descripteur unique n'est généralement suffisant pour représenter
l'ensemble des informations pertinentes contenues dans une image IRM. Cette limitation motive
les stratégies de \emph{fusion} de descripteurs étudiées dans ce travail. À titre
d'illustration, une étude comparative sur un jeu T1 de 2\,870 IRM montre qu'un descripteur
isolé plafonne autour de 85\,\% de précision (HOG avec un KNN de type \emph{subspace}), tandis
que la fusion GLCM~+~HOG~+~LBP porte la précision à 91,1\,\% avec un classifieur \emph{fine}
KNN~\cite{ref-pattanaik}. Ce gain confirme la complémentarité des descripteurs de texture et de forme.

\subsection{Pourquoi les étudier encore aujourd'hui~?}

Malgré la domination actuelle du Deep Learning, les approches basées sur les
caractéristiques manuelles conservent un intérêt scientifique et pratique important. Elles
demeurent particulièrement pertinentes lorsque les ressources de calcul sont limitées ou
lorsque la quantité de données disponibles est insuffisante pour entraîner efficacement des
réseaux profonds.

De plus, l'interprétabilité constitue aujourd'hui un enjeu majeur dans les applications
médicales. Les caractéristiques manuelles permettent de comprendre plus facilement les
éléments utilisés par les modèles pour prendre leurs décisions, ce qui favorise leur
acceptation dans les environnements cliniques. Enfin, elles ne s'opposent pas nécessairement
aux approches profondes~: de nombreux travaux récents fusionnent descripteurs profonds et
descripteurs manuels (par exemple CNN~+~GLCM ou CNN~+~HOG) afin de combiner sémantique de haut
niveau et sensibilité aux textures fines.

Dans ce contexte, ce projet s'inscrit dans la continuité de ces travaux en étudiant l'impact
de plusieurs familles de descripteurs manuels sur la classification multiclasses des tumeurs
cérébrales à partir d'images IRM, tout en évaluant leur complémentarité à travers différentes
stratégies de fusion et plusieurs algorithmes de Machine Learning.

% -----------------------------------------------------------------------
\section{Synthèse des travaux récents sur la classification des tumeurs cérébrales}

Au cours des dernières années, de nombreux travaux ont exploré l'utilisation des techniques
de Machine Learning et de Deep Learning pour la classification automatique des tumeurs
cérébrales à partir d'images IRM. Les performances rapportées dans la littérature ont connu
une progression importante grâce à l'utilisation de réseaux convolutifs profonds, de stratégies
de transfert d'apprentissage, ainsi que d'approches hybrides combinant extraction de
caractéristiques et modèles de classification.

Le Tableau~\ref{tab:etat-art} présente une synthèse de plusieurs travaux récents
représentatifs de l'état de l'art. Les méthodes recensées couvrent aussi bien les approches
traditionnelles basées sur des caractéristiques manuelles que les architectures profondes
modernes utilisant le transfert d'apprentissage ou la fusion de caractéristiques.

L'analyse de ces résultats montre que les architectures de Deep Learning atteignent
généralement des performances supérieures à 98\,\% de précision sur les jeux de données les
plus utilisés. Toutefois, ces performances reposent souvent sur des modèles complexes
nécessitant des ressources computationnelles importantes, des techniques d'augmentation de
données avancées et parfois des architectures hybrides difficiles à interpréter.

À l'inverse, les approches fondées sur des caractéristiques manuelles conservent plusieurs
avantages importants, notamment leur interprétabilité, leur faible coût de calcul et leur
capacité à fonctionner efficacement sur des ensembles de données de taille limitée. Cette
observation justifie l'intérêt de poursuivre l'étude des descripteurs classiques et de leurs
combinaisons dans un contexte médical où la transparence des décisions constitue un critère
essentiel.

Le présent projet s'inscrit dans cette perspective en proposant une étude approfondie de
plusieurs familles de caractéristiques manuelles (Statistiques, GLCM, LBP, DWT et HOG),
associées à un benchmark exhaustif de huit algorithmes de Machine Learning et à une
validation statistique rigoureuse des résultats obtenus.

\begin{table}[htbp]
  \centering
  \caption{Comparaison de travaux récents sur la classification des tumeurs cérébrales à
           partir d'images IRM.}
  \label{tab:etat-art}
  \scriptsize
  \setlength{\tabcolsep}{3.5pt}
  \renewcommand{\arraystretch}{1.2}
  \begin{tabularx}{\textwidth}{@{}
    >{\raggedright\arraybackslash}p{1.45cm}
    >{\raggedright\arraybackslash}X
    >{\raggedright\arraybackslash}p{2.35cm}
    >{\centering\arraybackslash}p{0.75cm}
    >{\centering\arraybackslash}p{0.85cm}
    >{\centering\arraybackslash}p{0.85cm}
    >{\centering\arraybackslash}p{0.85cm}
    >{\centering\arraybackslash}p{0.85cm}
  @{}}
    \toprule
    \thd{Réf.} & \thd{Méthode} & \thd{Dataset} & \thd{Cl.} &
    \thd{Acc.} & \thd{Prec.} & \thd{Rec.} & \thd{F1} \\
    \midrule
    Afzal et al.~\cite{ref-afzal}
      & ResNet18 + Transfer Learning + CART-ANOVA
      & Kaggle + Radiopaedia
      & 4 & \textbf{99.7} & -- & -- & 99.7 \\
    Ullah et al.~\cite{ref-ullahz}
      & VGG-16 / VGG-19 Transfer Learning
      & Brain Tumor MRI (Kaggle)
      & 2 & \textbf{99.2} & 99.0 & 99.0 & 99.0 \\
    Ullah et al.~\cite{ref-ullahn}
      & InceptionResNetV2 (Transfer Learning)
      & Glioma / Mening. / Pituit.
      & 3 & \textbf{98.9} & 98.3 & 99.8 & 99.0 \\
    Younis et al.~\cite{ref-younis}
      & VGG-16 Fine-Tuning
      & T1 MRI (253 images)
      & 2 & \textbf{98.5} & -- & 94.5 & 92.6 \\
    Shah et al.~\cite{ref-shah}
      & EfficientNet-B0 (Fine-Tuning)
      & Kaggle + BraTS2015 / TCIA
      & 2 & \textbf{98.9} & 98.9 & 98.8 & 98.8 \\
    Kundari et al.~\cite{ref-kundari}
      & DCNN + LBP/ICA + SVM
      & Glioma / Mening. / Pituit.
      & 3 & \textbf{99.0} & -- & 97.3 & -- \\
    Badža et al.~\cite{ref-badza}
      & CNN personnalisé (22 couches)
      & Figshare MRI Dataset
      & 3 & \textbf{96.6} & -- & -- & 97.5 \\
    Zahoor et al.~\cite{ref-zahoor}
      & BRAIN-RENet + HOG + SVM (HFF-BTC)
      & MRI Dataset
      & 3 & \textbf{99.2} & -- & -- & 99.1 \\
    Dixon et al.~\cite{ref-dixon}
      & CNN + ViT + Radiomics (GLCM/LBP) + MLP
      & Public + Local
      & 4 & \textbf{99.4} & -- & 99.5 & -- \\
    Pacal et al.~\cite{ref-pacal}
      & Swin Transformer + MLP résiduel
      & MRI Dataset
      & 4 & \textbf{99.9} & 99.9 & 99.8 & 99.9 \\
    \bottomrule
  \end{tabularx}
  \vspace{2pt}
  {\footnotesize\textit{Cl.} = nombre de classes ; métriques en \% (Accuracy, Precision, Recall, F1).}
\end{table}

\begin{notebox}[Note méthodologique~:]
Les valeurs reportées correspondent aux meilleures configurations décrites par les auteurs.
Certaines entrées sont issues d'une revue systématique récente~\cite{ref-review}, où les
protocoles d'évaluation (découpage des données, augmentation, validation croisée) ne sont pas
toujours homogènes. Les comparaisons directes entre lignes doivent donc être interprétées avec
prudence (voir Section~\ref{sec:analyse-critique}).
\end{notebox}

\subsection{Approches basées sur la fusion de caractéristiques manuelles}

Plusieurs travaux confirment spécifiquement l'apport de la fusion de descripteurs classiques,
qui constitue le c\oe{}ur de notre étude. Sur un jeu T1 de 2\,870 IRM à quatre classes, la
combinaison GLCM~+~HOG~+~LBP atteint 91,1\,\% de précision (fine KNN), contre au plus 85\,\%
pour un descripteur isolé, démontrant la complémentarité texture/forme~\cite{ref-pattanaik}. D'autres
études rapportent des combinaisons LBP~+~HOG associées à une réduction de dimension par ACP et
un classifieur SVM, ou encore l'ajout de \emph{features} d'interaction obtenues par produit
externe des vecteurs GLCM et LBP, atteignant des précisions de l'ordre de 96\,\% en
classification multiclasses~\cite{ref-basthikodi,ref-review}. Ces résultats situent les approches
purement manuelles dans une gamme de 85 à 96\,\%, inférieure aux modèles profonds les plus
récents, mais obtenue avec des modèles nettement plus légers, plus rapides à entraîner et plus
faciles à interpréter.

Le Tableau~\ref{tab:etat-art-ml} synthétise plusieurs travaux récents fondés exclusivement sur
des caractéristiques manuelles et des classifieurs classiques.

\begin{table}[htbp]
  \centering
  \caption{Travaux récents utilisant le Machine Learning traditionnel (caractéristiques
           manuelles et classifieurs classiques) pour la classification des tumeurs cérébrales
           à partir d'images IRM.}
  \label{tab:etat-art-ml}
  \scriptsize
  \setlength{\tabcolsep}{3.5pt}
  \renewcommand{\arraystretch}{1.2}
  \begin{tabularx}{\textwidth}{@{}
    >{\raggedright\arraybackslash}p{1.55cm}
    >{\raggedright\arraybackslash}X
    >{\raggedright\arraybackslash}p{2.35cm}
    >{\centering\arraybackslash}p{0.75cm}
    >{\centering\arraybackslash}p{0.85cm}
    >{\centering\arraybackslash}p{0.85cm}
    >{\centering\arraybackslash}p{0.85cm}
    >{\centering\arraybackslash}p{0.85cm}
  @{}}
    \toprule
    \thd{Réf.} & \thd{Méthode (descripteurs + classifieur)} & \thd{Dataset} &
    \thd{Cl.} & \thd{Acc.} & \thd{Prec.} & \thd{Rec.} & \thd{F1} \\
    \midrule
    Kumar et al.~\cite{ref-kumar}
      & Random Forest + Image Loading / HOG / LBP
      & Kaggle (2 datasets)
      & 2 & \textbf{99.0} & 99.0 & 97.0 & 98.0 \\
    Muthaiyan \& Malleswaran~\cite{ref-muthaiyan}
      & Bendlets (BCF, HPBC, HNBC) + Ensemble (kNN / NB / SVM)
      & REMBRANDT
      & 2 & \textbf{99.5} & -- & 99.0 & -- \\
    Basthikodi et al.~\cite{ref-basthikodi}
      & SVM + HOG + LBP + PCA
      & Kaggle (7023 img.)
      & 4 & \textbf{96.0} & 96.0 & 96.0 & 96.0 \\
    Pattanaik et al.~\cite{ref-pattanaik}
      & Fine KNN + fusion GLCM + HOG + LBP
      & Figshare (2870 img.)
      & 4 & \textbf{91.1} & -- & -- & -- \\
    \bottomrule
  \end{tabularx}
  \vspace{2pt}
  {\footnotesize\textit{Cl.} = nombre de classes ; métriques en \% (Accuracy, Precision, Recall, F1).}
\end{table}

\begin{notebox}[Note~:]
Pattanaik et al. rapportent une AUC de 0.96 pour la meilleure configuration (fusion
GLCM+HOG+LBP, \emph{fine} KNN)~; la précision d'un descripteur isolé y plafonne à 85\,\% (HOG,
\emph{subspace} KNN). Muthaiyan \& Malleswaran opèrent en deux étapes (normal/anormal puis
faible/haut risque)~: 99.5\,\% et 99\,\% respectivement. Kumar et al. obtiennent 99\,\% avec
Random Forest et le chargement d'image, SVM et régression logistique se situant entre 97 et
99\,\%.
\end{notebox}

\subsection{Analyse critique de l'état de l'art}
\label{sec:analyse-critique}

L'analyse des travaux récents met en évidence plusieurs tendances majeures.

Premièrement, les modèles de Deep Learning dominent actuellement la littérature grâce à leur
capacité à apprendre automatiquement des représentations complexes à partir des images IRM.
Les architectures fondées sur le transfert d'apprentissage (ResNet, VGG, EfficientNet,
InceptionResNet) et, plus récemment, les Transformers hiérarchiques (Swin) atteignent
régulièrement des taux de classification supérieurs à 99\,\%.

Deuxièmement, les approches hybrides associant réseaux profonds et classificateurs
traditionnels, ou fusionnant caractéristiques profondes et manuelles, égalent ou surpassent
souvent les modèles purement neuronaux. Les travaux combinant CNN~+~SVM, BRAIN-RENet~+~HOG~+~SVM
ou encore CNN~+~ViT~+~radiomique montrent que les informations extraites à différents niveaux
sont complémentaires.

Troisièmement, la majorité des études se concentrent sur l'amélioration des performances
prédictives, tandis que peu abordent en profondeur l'interprétabilité et la validation
statistique des résultats. Dans un contexte médical, cette limitation constitue un obstacle
important à l'adoption clinique.

Enfin, et surtout, les performances très élevées rapportées (souvent supérieures à 99\,\%)
doivent être interprétées avec prudence. Une revue récente souligne notamment qu'une précision
de 100\,\% a été obtenue par un CNN entraîné sur seulement 253 images IRM augmentées,
illustrant comment des jeux de données de très petite taille peuvent gonfler artificiellement
les estimations de performance~\cite{ref-review}. À cela s'ajoutent l'hétérogénéité des
protocoles d'évaluation, les fuites potentielles entre patients lors des découpages
image-par-image, et le manque de standardisation des jeux de données, qui limitent la
comparabilité directe des résultats et leur généralisation aux données cliniques réelles.

C'est dans cette optique que s'inscrit notre travail, qui vise à étudier systématiquement
l'apport de plusieurs familles de descripteurs manuels, leur complémentarité et leur impact
sur différents algorithmes de Machine Learning, à travers une méthodologie rigoureuse de
benchmarking et de validation statistique, privilégiant la transparence et la reproductibilité
plutôt que la seule maximisation de la précision.
